\documentclass[aspectratio=169,11pt]{beamer}
\usepackage{fontspec,polyglossia}
\setdefaultlanguage{russian}\setotherlanguage{english}
\setsansfont{DejaVu Sans}\setmainfont{DejaVu Sans}\setmonofont{DejaVu Sans Mono}
\newfontfamily\cyrillicfont{DejaVu Sans}\newfontfamily\cyrillicfontsf{DejaVu Sans}
\usepackage{amsmath,amssymb,booktabs,graphicx}
\definecolor{navy}{HTML}{102A43}\definecolor{teal}{HTML}{087E8B}
\setbeamercolor{structure}{fg=teal}\setbeamercolor{frametitle}{fg=navy}
\setbeamercolor{normal text}{fg=navy,bg=white}
\setbeamertemplate{navigation symbols}{}\setbeamertemplate{footline}[frame number]
\setbeamersize{text margin left=9mm,text margin right=9mm}
\title{Уравнения Максвелла}
\subtitle{EM scattering and tag-reader coupling\\Символьная модель, 3D-динамика и идентификация}
\author{Лабораторная работа · тема 3}
\date{29 октября 2026}
\newcommand{\fig}[2][.73]{\centering\includegraphics[height=#1\textheight,width=.98\textwidth,keepaspectratio]{../results/#2.png}}
\begin{document}
\begin{frame}\titlepage\end{frame}

\begin{frame}{Задача: от уравнений к проверяемой модели}
\begin{enumerate}
\item Сравнить символьные представления и векторные операторы.
\item Сохранить уравнения как словарь, граф и матрицу смежности.
\item Получить решение Максвелла в 3D со временем.
\item Проверить аппроксимацию и восстановить параметры ODE.
\item Подготовить интерфейс физической модели с поправкой.
\end{enumerate}
\vspace{3mm}
Три связанных учебных блока: \textbf{сфера → поле → катушки}. Это разные режимы одной теории; оптические данные не переносятся в RFID.
\end{frame}

\begin{frame}{Основание: журнальные статьи и открытые данные}
\small
\begin{tabular}{p{.32\textwidth}p{.59\textwidth}}
\toprule Источник & Роль в работе\\\midrule
Sharif et al., Sci. Rep., 2023 & Reader/tag как связанные цепи; геометрия статьи не воспроизводится.\\[2mm]
Abdelrahman et al., Sci. Rep., 2017 & Рассеяние на диэлектрических сферах, коэффициенты Ми.\\[2mm]
Malitson, 1965 & Дисперсия кварцевого стекла, 20°C.\\[2mm]
Daimon \& Masumura, 2007 & Дисперсия воды, 20°C.\\\bottomrule
\end{tabular}
\vspace{3mm}
Два независимых исследования материалов доставлены по HTTP через RefractiveIndex.INFO. Это опубликованные модели дисперсии; внешних измеренных 3D-полей здесь нет.
\end{frame}

\begin{frame}{Символьная модель в $(x,y,z,t)$}
\[
\epsilon\partial_t\mathbf E+\sigma\mathbf E+\mathbf J-\nabla\times\mathbf H=0,
\qquad \mu\partial_t\mathbf H+\nabla\times\mathbf E=0.
\]
\[\nabla\cdot\mathbf D=\rho,\qquad\nabla\cdot\mathbf B=0.\]
\begin{itemize}
\item SymPy Matrix: компоненты и производные видны явно.
\item sympy.vector: геометрическая запись curl, divergence.
\item SymEngine: независимый алгебраический движок.
\end{itemize}
Константные $\epsilon,\mu$ в этой записи; в численной неоднородной задаче интегрируется $\mathbf D=\epsilon(\mathbf x)\mathbf E$.
\end{frame}

\begin{frame}{Операторы: проверяем смысл и тождества}
\[\mathbf F=(yz,x^2z,xy^2),\quad f=x^2y+\sin z.\]
\[\nabla\times\mathbf F=(2xy-x^2,\ y-y^2,\ 2xz-z).\]
\begin{itemize}
\item grad, curl, div; скалярное и векторное произведение.
\item Матрица Якоби и покомпонентный лапласиан.
\item Проверки: $\nabla\times\nabla f=0$, $\nabla\cdot(\nabla\times\mathbf F)=0$.
\end{itemize}
Дивергенция возвращает скаляр. Совпадение представлений проверяется символически; единичные замеры времени не дают рейтинга скорости библиотек.
\end{frame}

\begin{frame}{Уравнение как граф вычисления}
\fig[.67]{expression_graph}
\small Ребро: родитель → операнд. Вершины — вхождения AST, а не точки пространственной сетки. Экспорт: JSON, GraphML, CSV.
\end{frame}

\begin{frame}{Матрица смежности хранит структуру выражения}
\begin{columns}
\column{.53\textwidth}\fig[.68]{adjacency}
\column{.44\textwidth}
$A_{ij}=1$, если узел $j$ — операнд узла $i$.
\vspace{4mm}
Порядок ID и подписи сохранены в отдельном JSON.
\vspace{4mm}
Словарь уравнений хранит \texttt{srepr}; обратное преобразование проверено тестом.
\end{columns}
\end{frame}

\begin{frame}{Рассеяние: Ми и предел Рэлея}
\[Q_{\rm sca}=\frac{2}{x^2}\sum_{n=1}^{N}(2n+1)(|a_n|^2+|b_n|^2),\quad x=2\pi a/\lambda.\]
\[Q_R=\frac83x^4\left|\frac{m^2-1}{m^2+2}\right|^2,\qquad x\ll1.\]
\begin{itemize}
\item Сфера в вакууме, $m=n_{\rm sphere}$, $\mu_r=1$.
\item $a=25$ nm, $\lambda=0.45\ldots1.0$ μm; прозрачные материалы.
\item Реальное поле: $\operatorname{Re}[\mathbf E(\mathbf r)e^{-i\omega t}]$.
\end{itemize}
Ми — независимый аналитический эталон для Рэлея. Периодический IVP далее имеет другие границы и начальные условия.
\end{frame}

\begin{frame}{Два материала: изменение дисперсии и рассеяния}
\fig[.66]{materials_scattering}
\small Максимальная относительная ошибка Рэлея на этом диапазоне: кварцевое стекло 0.36\%, вода 0.93\%. Это ошибка приближения относительно Ми, а не относительно измеренного поля.
\end{frame}

\begin{frame}{Рэлей имеет ограниченную область применимости}
\fig[.68]{rayleigh_error}
\small Ошибка $|Q_R-Q_M|/Q_M$, $m=1.46$. На малых $x$ приближение сходится к Ми; вне этого предела оно требует проверки.
\end{frame}

\begin{frame}{Численная задача: Maxwell 3D + время}
\[\dot{\mathbf D}=\nabla_h\times\mathbf H,\qquad
\dot{\mathbf H}=-\nabla_h\times(\mathbf D/\epsilon(\mathbf x)).\]
\small
\begin{itemize}
\item $[0,1)^3$, периодические границы; нормированные $c=\epsilon_0=\mu_0=1$.
\item Сглаженная сфера: $\epsilon=1+1.1\exp[-(r/0.14)^8]$.
\item Начальные $D_y=H_z=\cos(2\pi x)$; остальные компоненты нулевые.
\item Центральные разности $O(h^2)$ и RK4; 24³ узла, 160 шагов, $t_{\rm end}=0.65$.
\end{itemize}
\vspace{2mm}
Храним все компоненты E,H во времени. Нет PML и условия излучения: это периодическая начальная задача, а не модель изолированной антенны.
\end{frame}

\begin{frame}{Поле меняется в трёх пространственных координатах}
\fig[.65]{maxwell_slices}
\small Разность с однородной задачей показывает возмущение IVP. Начальные D,H одинаковы; E внутри включения различаются из-за ε. Это не far-field сечение рассеяния.
\end{frame}

\begin{frame}{Векторное возмущение поля в 3D}
\fig[.75]{maxwell_3d}
\end{frame}

\begin{frame}{Контроль дивергенции и энергии}
\fig[.65]{conservation}
\small $\max|\operatorname{div}D|\approx3.4\cdot10^{-14}$, $\max|\operatorname{div}H|\approx2.9\cdot10^{-14}$. Максимальный относительный дрейф энергии $\approx1.7\cdot10^{-9}$.
\end{frame}

\begin{frame}{Проверка порядка на точном решении}
\begin{columns}
\column{.60\textwidth}\fig[.68]{convergence}
\column{.36\textwidth}\small
Эталон: $E_y=H_z=\cos[2\pi(x-t)]$ в однородной среде.
\vspace{4mm}
\begin{tabular}{rr}\toprule $n$ & RMSE\\\midrule 12 & 0.12997\\24 & 0.03287\\48 & 0.00824\\\bottomrule\end{tabular}
\vspace{4mm}
Порядки: 1.98 и 2.00.
\end{columns}
\end{frame}

\begin{frame}{От 3D-геометрии к связи reader/tag}
\begin{columns}
\column{.54\textwidth}\fig[.65]{coils_3d}
\column{.43\textwidth}\small
\[M=\frac{\mu_0}{4\pi}\oint\!\oint\frac{d\boldsymbol\ell_1\cdot d\boldsymbol\ell_2}{|\mathbf r_1-\mathbf r_2|}.\]
Один виток; радиусы 40 и 15 mm; расстояние 40 mm. Квадратура: 240 точек на петлю.
\vspace{3mm}
$f=13.56$ MHz, $\lambda\approx22$ m. Геометрия намного меньше λ: квазистатическая модель допустима.
\end{columns}
\end{frame}

\begin{frame}{Связь зависит от расстояния и ориентации}
\fig[.70]{coupling_map}
\small Проверены взаимность $M_{12}=M_{21}$, сходимость квадратуры и нулевая связь для ортогональных соосных петель.
\end{frame}

\begin{frame}{Редуцированная модель динамики контуров}
\[\underbrace{\begin{pmatrix}L_1&M\\M&L_2\end{pmatrix}}_{\mathbf L}\dot{\mathbf i}
+\mathbf R\mathbf i+\mathbf C^{-1}\mathbf q=(V_s,0)^T,\qquad\dot{\mathbf q}=\mathbf i.\]
\small
\begin{itemize}
\item $L_1=1$ μH, $L_2=0.5$ μH, $R_1=2$ Ω, $R_2=4$ Ω.
\item $C_j=[(2\pi f)^2L_j]^{-1}$; нулевые начальные токи и заряды.
\item $V_s=\cos(2\pi ft)$ V, $f=13.56$ MHz; решатель DOP853.
\item Пассивность: $M^2<L_1L_2$.
\end{itemize}
Вектор подбора $\theta=(M\,[\mathrm{nH}],R_2\,[\Omega])$. Начальная точка (12,6), границы [0.1,100] и [0.5,15].
\end{frame}

\begin{frame}{Подбор параметров по синтетическому сигналу}
\fig[.62]{parameter_fit}
\small 501 точка за 2 μs. Gaussian шум: 1\% peak current; seed фиксирован. Каждая третья точка исключена из least squares. Истинные данные и оценка получены с разными точностями ODE.
\end{frame}

\begin{frame}{Результат идентификации и его границы}
\centering
\begin{tabular}{lrr}\toprule Параметр & Истина & Оценка\\\midrule
$M$, nH & 3.7673 & 3.8262\\
$R_2$, Ω & 4.0000 & 4.0730\\\bottomrule
\end{tabular}
\vspace{5mm}
\begin{itemize}
\item Held-out RMSE относительно чистого сигнала: $2.74\cdot10^{-5}$ A.
\item Стандартное отклонение шума: $3.24\cdot10^{-4}$ A.
\item Якобиан имеет ранг 2 в найденной точке.
\end{itemize}
\vspace{2mm}
\small Это проверка на отложенных точках той же синтетической траектории. Она не доказывает перенос на новую антенну или реальные измерения.
\end{frame}

\begin{frame}{Структура для смешанной модели}
\[M_{\rm hybrid}=M_{\rm geometry}(d,\alpha)+\Delta M_{\phi}(d,\alpha,\ldots).\]
\begin{enumerate}
\item Физический блок вычисляет M по геометрии.
\item Пользовательская функция добавляет поправку в henry.
\item Проверяется положительность L-матрицы.
\item ODE-солвер строит предсказание тока.
\end{enumerate}
\vspace{3mm}
\texttt{HybridModel(correction=callable)} реализован. По умолчанию $\Delta M=0$. Нейросеть не обучалась; для неё нужны независимые измерения.
\end{frame}


\begin{frame}{Что удалось проверить}
\begin{itemize}
\item Символьные уравнения можно переводить в словарь и граф без потери структуры.
\item Сохранение ограничений и сходимость дают количественную проверку Maxwell IVP.
\item Приближение Рэлея проверяется на независимом решении Ми.
\item Параметры M и R₂ восстанавливаются в контролируемом synthetic experiment.
\item Геометрия, физика и поправка соединены через явный интерфейс.
\end{itemize}
\vspace{4mm}
Следующий эксперимент: добавить реальные измерения, шумовые возмущения и проверку на новых положениях метки.
\end{frame}

\begin{frame}{Источники}
\scriptsize
\begin{enumerate}
\item Sharif et al. Scientific Reports 13,14017 (2023). DOI: 10.1038/s41598-023-39825-9.
\item Abdelrahman et al. Scientific Reports 7,14762 (2017). DOI: 10.1038/s41598-017-15192-0.
\item Tzarouchis, Sihvola. Applied Sciences 8,184 (2018). DOI: 10.3390/app8020184.
\item Malitson. JOSA 55,1205–1208 (1965). DOI: 10.1364/JOSA.55.001205.
\item Daimon, Masumura. Applied Optics 46,3811–3820 (2007). DOI: 10.1364/AO.46.003811.
\item Polyanskiy. Scientific Data 11,94 (2024). DOI: 10.1038/s41597-023-02898-2.
\end{enumerate}
\vspace{4mm}
Полные HTTP-адреса, роль формул, документация библиотек: \texttt{docs/SOURCES.md}. Фактические численные результаты: \texttt{results/metrics.json}.
\end{frame}
\end{document}
